
\section{Mathematical Derivation}

\subsection{Integer Kernel}

For a primitive Miller index $\mathbf h=(h,k,l)$, define the integer kernel
\[
K(\mathbf h)=\{\mathbf x\in\mathbb Z^3:\mathbf h\cdot\mathbf x=0\}.
\]
Since $\mathbf h$ is primitive, $K(\mathbf h)$ is a rank-two sublattice of
$\mathbb Z^3$.

\subsection{Existence of a Bézout Witness}

By Bézout's identity there exists $\mathbf w\in\mathbb Z^3$ satisfying
\[
\mathbf h\cdot\mathbf w=1.
\]

\begin{theorem}
There exists an orientation-preserving unimodular matrix
$S\in SL(3,\mathbb Z)$ such that
\[
\mathbf hS=(0,0,1).
\]
\end{theorem}

\begin{proof}[Proof sketch]
Choose a primitive basis of $K(\mathbf h)$ together with a Bézout witness
$\mathbf w$. The resulting matrix is unimodular. A final permutation and,
if necessary, sign change produce an orientation-preserving basis.
\end{proof}

\section{Algorithm}

\begin{enumerate}
\item Reduce the Miller index to primitive form.
\item Construct Bézout coefficients.
\item Build a unimodular basis.
\item Permute basis vectors.
\item Repair orientation.
\item Verify invariants.
\end{enumerate}
